\documentclass[10pt,a4paper]{amsart}
\usepackage[margin=19mm]{geometry}
\usepackage{amsmath,amssymb,mathtools}
\usepackage[scr=rsfs,cal=euler]{mathalfa}
\usepackage{xfrac}
\usepackage[unicode,hidelinks,bookmarks=false]{hyperref}
\newcommand{\ie}{i.e.\ }
\newcommand{\cf}{cf.\ }
\newcommand{\resp}{respectively\ }
\newcommand{\Subsection}{\subsection}
\providecommand{\nobreakdash}{\nobreak\hbox{-}}
\setlength{\emergencystretch}{2em}
\multlinegap0pt
\usepackage{amssymb}
\usepackage{cancel}
\usepackage[all]{xy}
\usepackage{xcolor}
\usepackage{tikz}
\usepackage{csquotes}
\usepackage{bbm}
\usepackage{mathtools}
\usepackage{caption}
\usepackage{booktabs}
\newtheorem{thm}{}
\title{Components of discriminants for systems of equations and irreducibility of determinants}
\author{Vladislav Pokidkin}
\date{Source: arXiv:2501.15832v2 (CC BY 4.0)}
\begin{document}
\maketitle

\noindent\emph{Source and licence.} Components of discriminants for systems of equations and irreducibility of determinants. Version arXiv:2501.15832v2 (CC BY 4.0), distributed by arXiv under the Creative Commons Attribution 4.0 International licence, http://creativecommons.org/licenses/by/4.0/, as recorded in the arXiv licence metadata for this exact version and shown on the abstract page https://arxiv.org/abs/2501.15832v2. Full licence notice: \texttt{licenses/arxiv-2501.15832-CC-BY-4.0.txt}.\par\smallskip

\noindent\textit{Editorial scope.} Counted unit: the article's thm Conjecture (source0.tex lines 599--603). The article's complete original proof of it is delivered with it (source0.tex lines 605--669, one \texttt{proof} environment of 65 lines). No mathematics is written, repaired or re-derived: the statement and the proof are the article's own lines, in the article's own notation, and no retained line is substituted or re-flowed.  Retained uncounted as the source-local context the proof needs: lines 555--560.  Nothing was omitted from any retained span, so the package carries no omission record.  No retained line contains a citation, so the wrapper carries no bibliography and no bibliography entry.  No retained line contains a figure environment, an image-inclusion command or drawing code, so no image is delivered and the package declares no asset dependency.  Editorial formatting only, in addition: an AMS \texttt{amsart} preamble that restates the article's own declarations and macros used by the retained lines, namely the environments \texttt{thm} on separate counters, together with the standard packages those lines need.  The retained environments print in this document as Theorem 1 in order of appearance, whereas the article prints the same items with the numbering of its own full document.  The complete native source of the article as distributed in the arXiv e-print for this exact version is bundled unchanged as \texttt{sources/172223/source0.tex}.
\begin{thm}
\label{Theorem. Intersection of the determinant with a subspace}For
an irreducible BK-tuple $(L_{1},...,L_{n})$ of vector subspaces,
the intersection of the determinant hypersurface $Det$ with the subspace
$L$ is a variety.
\end{thm}

\begin{thm}
\label{Theorem. Esterov's Conjecture}For an irreducible BK-tuple
$\mathscr{A}$, the discriminant $D_{\mathscr{A}}$ is irreducible
in $\mathbb{K}_{\mathscr{A}}$.
\end{thm}

\begin{proof}
A tuple $\mathscr{A}$ defines a vector subspace of matrices $Mat_{\mathscr{A}}=\underset{A\in\mathscr{A}}{\oplus}(\langle A\rangle\otimes\mathbb{K})$
in the space of matrices $Mat_{dim\langle\mathscr{A}\rangle,\mathfrak{c}(\mathscr{A})}$.
For the irreducible BK-tuple $\mathscr{A}$, the intersection $Y$
of the determinant hypersurface $Det$ with the subspace $Mat_{\mathscr{A}}$
is a variety by Theorem \ref{Theorem. Intersection of the determinant with a subspace}.

With a finite set $A$ and a tuple $\mathscr{A}$, we associate matrices
$\tau_{A}=(a,\;a\in A)$ and $\tau_{\mathscr{A}}=\underset{A\in\mathscr{A}}{\oplus}\tau_{A}$,
with integer coefficients and ranks $rk\,\tau_{A}=dim\,\langle A\rangle,$
$rk\,\tau_{\mathscr{A}}=dim\,Mat_{\mathscr{A}}$. The matrix $\tau_{\mathscr{A}}$
defines the linear map from $\mathbb{K}_{\mathscr{A}}$ to $Mat_{\mathscr{A}}$
by the formula $\tau_{\mathscr{A}}(\lambda_{\mathscr{A}})=(\tau_{A}(\lambda_{A}),\;A\in\mathscr{A})$,
where $\tau_{A}(\lambda_{A})=\underset{a\in A}{\sum}\lambda_{a}\,a$.
This means that the preimage $K=\tau_{\mathscr{A}}^{-1}(Y)$ is a
variety in $\mathbb{K}_{\mathscr{A}}$ as a trivial vector bundle
of rank $dim\,\mathbb{K}_{\mathscr{A}}-dim\,Mat_{\mathscr{A}}$.

We can construct the following sequence of maps:
\[
\begin{array}{ccccccc}
T(\mathscr{A})\times\mathbb{K}_{\mathscr{A}} & \overset{\sigma_{\mathscr{A}}}{\twoheadrightarrow} & \mathbb{K}_{\mathscr{A}} & \overset{\tau_{\mathscr{A}}}{\rightarrow} & Mat_{\mathscr{A}} & \overset{det}{\rightarrow} & \mathbb{K},\end{array}
\]
where $\sigma_{\mathscr{A}}(x,c)=(c_{a}x^{a})_{a\in A\in\mathscr{A}}$
is the torus action. Notice that the preimage $W=\sigma_{\mathscr{A}}^{-1}(K)$
is a variety that is isomorphic to $T(\mathscr{A})\times K$. Indeed,
the isomorphism is defined by the regular maps $W\stackrel[\iota]{\sigma}{\rightleftarrows}T(\mathscr{A})\times K$
such that $\iota(x,c_{a})=(x,\frac{c_{a}}{x^{a}})$ and $\sigma(x,c_{a})=(x,c_{a}x^{a})$
for every element $a\in A\in\mathscr{A}.$

The sequence of maps defines the equation for a root $x$ of a polynomial
system $c=(c_{a})_{a\in A\in\mathscr{A}}$ to be a singular point:
$det\circ\tau_{\mathscr{A}}\circ\sigma_{\mathscr{A}}(c,x)=det\,||\underset{a\in A}{\sum}c_{a}x^{a}a||_{A\in\mathscr{A}}=0.$
This polynomial is irreducible because $W$ is a variety.

Consider the vector subspaces $\Pi_{A}$ in $\mathbb{K}_{\mathscr{A}}$
defined by the linear equations $\underset{a\in A}{\sum}c_{a}=0$,
and their intersection $\Pi=\underset{A\in\mathscr{A}}{\cap}\Pi_{A}$.
Denote by $c_{A}$ the coefficient corresponding to zero in a finite
set $A$. In each linear equation, we can express the variable $c_{A}=-\underset{a\in A\backslash\{0\}}{\sum}c_{a}.$
Notice that the determinant $det\,||\underset{a\in A}{\sum}c_{a}a||_{A\in\mathscr{A}}=0$
does not depend on the coefficients $c_{A},$ $A\in\mathscr{A}$.
Then the intersection of the variety $K$ with the subspace $\Pi$
is a variety $X=K\cap\Pi$ defined by the same determinant. To see
this, we ensure that the coordinate algebra $\mathbb{K}[X]$ is an
integral domain:
\[
\mathbb{K}[X]=\frac{\mathbb{K}[c_{a},\,a\in A\in\mathscr{A}]}{\langle\underset{a\in A}{\sum}c_{a},\;A\in\mathscr{A};\;det\,||\underset{a\in A}{\sum}c_{a}a||_{A\in\mathscr{A}}\rangle}=\frac{\mathbb{K}[c_{a},\,a\in A\backslash\{0\},\,A\in\mathscr{A}]}{\langle det\,||\underset{a\in A}{\sum}c_{a}a||_{A\in\mathscr{A}}\rangle}.
\]

Thus the preimage $Z=\sigma_{\mathscr{A}}^{-1}(X)$ is a variety,
isomorphic to the direct product $T(\mathscr{A})\times X$ by the
same isomorphisms $\iota$ and $\sigma$. Notice that each preimage
$\sigma_{\mathscr{A}}^{-1}(\Pi_{A})$ equals the zero locus of the
equation $f_{A}(x)=0$ in $T(\mathscr{A})\times\mathbb{K}_{\mathscr{A}}$.
Hence the variety $Z$ projects to the discriminant $D_{\mathscr{A}}$
by the map $T(\mathscr{A})\times\mathbb{K}_{\mathscr{A}}\overset{\pi}{\rightarrow}\mathbb{K}_{\mathscr{A}}$,
$\pi(x,c)=(c)$:
\begin{align*}
Z=\sigma_{\mathscr{A}}^{-1}(X) & =\sigma_{\mathscr{A}}^{-1}(K\cap\Pi)=\sigma_{\mathscr{A}}^{-1}(K)\cap\sigma_{\mathscr{A}}^{-1}(\Pi)=W\underset{A\in\mathscr{A}}{\cap}\sigma_{\mathscr{A}}^{-1}(\Pi_{A})=\\
 & =V(det\,||\underset{a\in A}{\sum}c_{a}x^{a}a||_{A\in\mathscr{A}})\cap\underset{A\in\mathscr{A}}{\cap}V(f_{A}(x)).
\end{align*}
Therefore, the discriminant $D_{\mathscr{A}}$ is a variety as a projection
of a variety.
\end{proof}

\end{document}
